\documentclass[11pt,twocolumn]{article}

% ── Packages ────────────────────────────────────────────────────
\usepackage[margin=1in]{geometry}
\usepackage{times}
\usepackage{microtype}
\usepackage{amsmath,amssymb}
\usepackage{graphicx}
\usepackage{booktabs}
\usepackage{tabularx}
\usepackage{multirow}
\usepackage{array}
\usepackage{xcolor}
\usepackage{hyperref}
\usepackage{url}
\usepackage{enumitem}
\usepackage{fancyhdr}
\usepackage{abstract}
\usepackage{titlesec}
\usepackage{caption}
\usepackage{float}
\usepackage{balance}
\usepackage{natbib}
\usepackage{mdframed}

% ── Colours ─────────────────────────────────────────────────────
\definecolor{teal}{RGB}{14,124,123}
\definecolor{navy}{RGB}{27,42,74}
\definecolor{amber}{RGB}{232,168,56}
\definecolor{lightgray}{RGB}{240,244,248}
\definecolor{darkgray}{RGB}{51,65,85}

% ── Hyperref setup ───────────────────────────────────────────────
\hypersetup{
  colorlinks=true,
  linkcolor=navy,
  citecolor=teal,
  urlcolor=teal,
  pdftitle={AI FinOps: Why Cloud Financial Governance Fails for Intelligent Systems},
  pdfauthor={Sathiyan Kutty}
}

% ── Section heading style ────────────────────────────────────────
\titleformat{\section}{\large\bfseries\color{navy}}{{\thesection}}{1em}{}[\color{teal}\titlerule]
\titleformat{\subsection}{\normalsize\bfseries\color{darkgray}}{\thesubsection}{1em}{}
\titlespacing{\section}{0pt}{10pt}{4pt}
\titlespacing{\subsection}{0pt}{6pt}{2pt}

% ── Callout box (mdframed) ───────────────────────────────────────
\newmdenv[
  backgroundcolor=lightgray,
  linecolor=teal,
  linewidth=1.5pt,
  leftline=true, rightline=false, topline=false, bottomline=false,
  innerleftmargin=8pt, innerrightmargin=8pt,
  innertopmargin=6pt, innerbottommargin=6pt,
  skipabove=6pt, skipbelow=6pt
]{callout}

% ── Header / footer ─────────────────────────────────────────────
\pagestyle{fancy}
\fancyhf{}
\fancyhead[L]{\small\textcolor{darkgray}{Kutty \textbar\ AI FinOps}}
\fancyhead[R]{\small\textcolor{darkgray}{FinOps Foundation White Paper \textbar\ 2026}}
\fancyfoot[C]{\small\thepage}
\renewcommand{\headrulewidth}{0.4pt}

% ── Abstract style ───────────────────────────────────────────────
\renewcommand{\abstractnamefont}{\large\bfseries\color{navy}}
\renewcommand{\abstracttextfont}{\small}
\setlength{\absleftindent}{0pt}
\setlength{\absrightindent}{0pt}

% ── Caption style ────────────────────────────────────────────────
\captionsetup{font=small, labelfont={bf,color=navy}, textfont=it}

% ── Table column types ───────────────────────────────────────────
\newcolumntype{L}[1]{>{\raggedright\arraybackslash}p{#1}}
\newcolumntype{C}[1]{>{\centering\arraybackslash}p{#1}}

% ─────────────────────────────────────────────────────────────────
\begin{document}

% ── Title block ──────────────────────────────────────────────────
\twocolumn[{%
\begin{@twocolumnfalse}

\vspace*{4pt}
{\color{teal}\rule{\linewidth}{2pt}}
\vspace{6pt}

{\noindent\LARGE\bfseries\color{navy}
AI FinOps: Why Cloud Financial Governance Fails\\[4pt]
for Intelligent Systems --- and a Framework\\[4pt]
for What Replaces It}

\vspace{10pt}

{\noindent\large\bfseries Sathiyan Kutty}\\[2pt]
{\noindent\normalsize Chief AI Officer, EMIDS Healthcare IT\\
\href{https://www.linkedin.com/in/skutty/}{linkedin.com/in/skutty}
\hfill
\textit{Submitted to the FinOps Foundation, March 2026}}

\vspace{8pt}
{\color{teal}\rule{\linewidth}{0.6pt}}
\vspace{10pt}

% ── Abstract ────────────────────────────────────────────────────
\begin{onecolabstract}
\noindent
Enterprise AI spending is growing at an unprecedented rate, yet the financial governance
frameworks organisations rely on were designed for a fundamentally different cost structure.
Cloud FinOps --- the dominant discipline for managing cloud infrastructure expenditure ---
operates on compute-hour billing, resource-level attribution, and predictable linear
scaling. AI systems, by contrast, are denominated in tokens, attributed at the
request level, and exhibit nonlinear cost dynamics driven by reasoning chains and
autonomous agentic loops that can multiply per-session consumption by one to three
orders of magnitude. These are architectural differences, not implementation gaps,
and they render existing Cloud FinOps frameworks insufficient for governing AI spend.

This paper makes four original contributions. First, it provides a formal definition
of \textit{AI FinOps} as a distinct organisational discipline. Second, it introduces
an eight-dimension AI Cost Taxonomy that identifies every measurable cost category
in a production AI deployment and maps each to an attribution mechanism and
optimisation lever. Third, it proposes a three-level AI FinOps Maturity Model
--- Reactive, Managed, and Optimised --- with specific, assessable capabilities
at each level. Fourth, it defines five governance questions that constitute a
minimum viable baseline assessment for any organisation deploying AI at scale.
Empirical grounding draws on published benchmarks, Gartner and Uptime Institute
research, and anonymised patterns observed across enterprise healthcare AI deployments.
The paper argues that AI FinOps is not optional overhead: organisations that fail to
build this capability face a projected 50\% cost penalty relative to governed peers
by 2030, while those that master it free the capital required to build the proprietary
AI assets that create durable competitive advantage.

\smallskip
\noindent\textbf{Keywords:} AI FinOps, token economics, cloud financial governance,
agentic AI, enterprise AI cost management, maturity model, AI governance
\end{onecolabstract}

\vspace{10pt}
{\color{teal}\rule{\linewidth}{0.6pt}}
\vspace{12pt}

\end{@twocolumnfalse}
}]

% ─────────────────────────────────────────────────────────────────
\section{Introduction}
\label{sec:intro}

Consider the following scenario, representative of a pattern observed across
enterprise AI deployments: a software engineering team deploys an AI-powered
research assistant on a Friday afternoon. By Monday morning, the organisation
has received a cloud invoice for \$250 in token consumption --- roughly 10$\times$
the team's weekly expectation. No one authorised the spend. No alert fired. No
governance control existed to prevent it. The system had simply been running,
autonomously generating reasoning chains and tool calls through a weekend with
no human in the loop and no budget constraint in place \citep{anthropic2025}.

This scenario is not an edge case. It is an archetype. As enterprise AI deployment
scales from isolated pilots to production infrastructure, the gap between the rate
of AI cost accumulation and the rate of governance capability development has
become a primary financial risk. Gartner projects that organisations failing to
optimise their AI compute environment will pay a 50\% cost penalty relative to
governed peers by 2030 \citep{gartner2025forecast}. Total global enterprise AI
spending reached approximately \$1.5 trillion in 2025, with projections of \$2.5
trillion in 2026 \citep{gartner2025spending}. The stakes of ungoverned consumption
at this scale are material.

The dominant financial governance framework for cloud infrastructure --- Cloud
FinOps, as formalised by the FinOps Foundation \citep{finopsfoundation2023} ---
was designed to address this kind of governance gap for compute resources. Its
lifecycle of Inform, Optimise, and Operate has proven effective for managing
virtual machine fleets, storage systems, and API-based services billed by the
compute hour. The FinOps Foundation's maturity model, persona framework, and
capability domains represent genuine institutional knowledge built over a decade
of enterprise cloud cost management practice.

The central argument of this paper is that Cloud FinOps, despite its maturity,
is \textit{architecturally insufficient} for AI systems. This is not a criticism
of Cloud FinOps --- it is an observation that the cost structure of AI is
different from the cost structure of cloud infrastructure in ways that are
fundamental, not superficial. Token-based billing, request-level attribution,
reasoning chain multiplication, and autonomous agentic dynamics introduce
cost variables that cloud billing APIs do not expose and that cloud cost
optimisation levers cannot address.

A new discipline is required. We define it as \textit{AI FinOps}: the
organisational capability comprising the instrumentation, attribution, governance,
and optimisation practices required to manage AI system costs at the token level,
connecting AI expenditure to business outcomes across model inference, agentic
workloads, and training operations.

This paper is structured as follows. Section~\ref{sec:background} reviews the
evolution of IT financial governance and the theoretical dynamics that make AI
cost management urgent. Section~\ref{sec:failure} demonstrates, across eight
dimensions, why Cloud FinOps fails for AI. Section~\ref{sec:framework} presents
the four original contributions of this paper. Section~\ref{sec:implementation}
describes a practitioner-validated 30-day implementation sprint. Section~\ref{sec:evidence}
provides empirical validation. Section~\ref{sec:discussion} discusses strategic
implications, and Section~\ref{sec:conclusion} concludes.

% ─────────────────────────────────────────────────────────────────
\section{Background and Related Work}
\label{sec:background}

\subsection{The Evolution of IT Financial Governance}

The governance of information technology costs has evolved through three
recognisable phases. The first phase --- IT cost accounting --- emerged from
ITIL service management frameworks in the 1990s, establishing that IT
expenditure should be attributed to business units and managed against budgets
\citep{itilv4}. The second phase --- cloud cost management --- emerged with
the mainstream adoption of hyperscaler infrastructure in the 2010s, when
on-demand billing models introduced the risk of unbounded consumption that
fixed-capacity data centre procurement had naturally constrained. The FinOps
Foundation formalised this phase, defining a shared vocabulary, a capability
model, and an organisational practice for cloud cost governance
\citep{finopsfoundation2023}.

The third phase is now underway: the governance of AI system costs. Like the
transition from data centre to cloud, this transition is not merely an extension
of the prior phase. It introduces a qualitatively different cost structure
--- one denominated in tokens rather than compute hours, exhibiting nonlinear
rather than linear scaling, and requiring attribution at the level of individual
model inference requests rather than provisioned resources.

\subsection{The AI Cost Structure: Tokens, Agents, and Nonlinearity}

The fundamental billing unit of AI API services is the token --- approximately
3--4 characters or 0.75 words in English \citep{openai2024}. Token pricing is
asymmetric: output tokens are priced at three to five times the rate of input
tokens due to the sequential, compute-intensive nature of autoregressive
generation. In 2026, prices range from \$0.05 per million tokens for nano-class
models to \$168 per million tokens for frontier reasoning models at peak output
rates --- a spread of three orders of magnitude \citep{introl2025}.

Two cost dynamics distinguish AI from prior cloud cost categories. The first
is \textit{reasoning chain multiplication}: when extended thinking or chain-of-thought
reasoning is enabled, models generate internal token streams before producing
visible output. These reasoning tokens are billed at output rates and can
multiply per-query cost by a factor of 6--10$\times$ relative to standard inference
\citep{anthropic2025extended}. The second is \textit{agentic loop compounding}:
autonomous AI agents executing multi-step tasks accumulate token costs across
iterations, tool calls, retrieval operations, and state management. Under
conservative assumptions, a 10-iteration agentic session may consume 50,000--80,000
tokens per session; under pathological conditions, consumption can exceed
1 million tokens \citep{introl2025}.

Neither dynamic is visible in standard cloud billing dashboards, and neither
is addressable through the reservation, savings plan, or right-sizing mechanisms
that constitute the Cloud FinOps optimisation toolkit.

\subsection{The Jevons Paradox in AI Compute}

The urgency of AI FinOps is compounded by a macroeconomic dynamic first
identified by William Stanley Jevons in 1865 in the context of coal and steam
engine efficiency \citep{jevons1865}. Jevons observed that improvements in
steam engine efficiency did not reduce total coal consumption; instead, lower
operating costs made steam power economically viable for a broader range of
applications, expanding total demand. This rebound effect --- now known as
the Jevons Paradox --- applies with precision to AI token economics.

Token prices have declined approximately 10$\times$ per year since 2022
\citep{introl2025}. This deflation has not produced declining enterprise AI
budgets. Instead, each generation of lower prices has enabled a new class of
AI use case --- ambient documentation, autonomous prior authorisation processing,
multi-agent research workflows --- that consumes tokens at volumes that absorb
and exceed the savings from unit cost reductions. Governance is therefore not
rendered unnecessary by falling prices; it is made \textit{more} necessary, because
the lower the unit cost, the faster consumption expands when ungoverned.

% ─────────────────────────────────────────────────────────────────
\section{Why Cloud FinOps Fails for AI}
\label{sec:failure}

The failure of Cloud FinOps for AI is not a failure of implementation but a
failure of architectural fit. Table~\ref{tab:comparison} documents eight
dimensions on which the cost structures of cloud infrastructure and AI systems
differ in ways that require different governance approaches.

\begin{table*}[t]
\centering
\caption{Cloud FinOps vs.\ AI FinOps: Architectural Comparison Across Eight Dimensions}
\label{tab:comparison}
\small
\begin{tabularx}{\textwidth}{L{2.2cm} L{4.6cm} L{4.6cm}}
\toprule
\textbf{Dimension} & \textbf{Cloud FinOps} & \textbf{AI FinOps (Required)} \\
\midrule
Cost unit &
  Compute hours, storage GB, API calls &
  Tokens (input, output, reasoning) --- invisible in standard billing dashboards \\
\addlinespace
Attribution model &
  Resource-level tagging at provisioning time &
  Request-level instrumentation in application middleware \\
\addlinespace
Anomaly detection &
  Spend threshold alerts on compute/storage &
  Token rate monitoring, reasoning chain length alerts, iteration count tracking \\
\addlinespace
Governance levers &
  Reserved instances, savings plans, right-sizing &
  Per-agent token budgets, iteration limits, reasoning token caps, routing policies \\
\addlinespace
Budget predictability &
  High --- VM hours and storage scale linearly &
  Low --- reasoning chains and agentic loops create nonlinear cost spikes \\
\addlinespace
Vendor landscape &
  1--3 hyperscalers; mature enterprise contracts &
  5--10 model providers + open-source deployments; rapidly evolving pricing \\
\addlinespace
Team ownership &
  Centralised FinOps team with cloud billing expertise &
  Distributed: engineering, data science, and finance share accountability \\
\addlinespace
Tooling maturity &
  Mature ecosystem (CloudHealth, Apptio, AWS Cost Explorer) &
  No dedicated AI FinOps tooling exists as of 2026 \\
\bottomrule
\end{tabularx}
\end{table*}

The most consequential of these differences is attribution. Cloud FinOps governance
depends on tagging resources --- virtual machines, buckets, databases --- at
the time of provisioning. Tags persist throughout the resource lifecycle and are
automatically captured in billing records. AI cost attribution cannot work this way:
the cost unit is not a provisioned resource but an ephemeral inference request.
Attribution requires application-layer instrumentation --- middleware that logs
token counts, model identifiers, task types, and session identifiers on every
call. Without this instrumentation, AI spend is a single aggregate line item,
opaque to the analysis that governance requires.

The second critical difference is the nature of anomalies. Cloud cost anomalies
are typically caused by provisioning errors or sudden traffic growth ---
events that manifest as elevated compute-hour consumption over minutes to hours.
AI cost anomalies can be caused by reasoning model activation on a simple task,
an agentic loop that fails to terminate, or a retrieval operation that injects
50,000 tokens of context where 500 would suffice. These anomalies can materialise
within a single session, at rates that are invisible to hourly billing checks.

\begin{callout}
\small\textbf{Observed pattern:} Across enterprise healthcare AI deployments,
the first AI FinOps assessment consistently reveals at least one application
consuming 3--5$\times$ its team's believed spend, typically due to uncontrolled
reasoning token usage or uncapped agentic sessions. The discovery predates
any month-end billing cycle by an average of 11 days when token-level
instrumentation is implemented.
\end{callout}

% ─────────────────────────────────────────────────────────────────
\section{The AI FinOps Framework}
\label{sec:framework}

This section presents the four original contributions of this paper.

\subsection{Definition of AI FinOps}
\label{sec:definition}

We define AI FinOps as follows:

\begin{mdframed}[backgroundcolor=lightgray, linecolor=teal, linewidth=1.5pt,
  innerleftmargin=10pt, innerrightmargin=10pt,
  innertopmargin=8pt, innerbottommargin=8pt]
\small\textit{\textbf{AI FinOps} is the organisational capability comprising the
instrumentation, attribution, governance, and optimisation practices required to
manage AI system costs at the token level, with accountability structures that
connect AI expenditure to measurable business outcomes across model inference,
agentic workloads, and training operations.}
\end{mdframed}

\vspace{4pt}
This definition has three structural components. \textit{Instrumentation and attribution}
refers to the technical infrastructure required to measure and assign token costs
at the request level --- the necessary precondition for any governance action.
\textit{Governance} refers to the controls --- budgets, limits, routing policies,
alerting thresholds --- that bound consumption within authorised parameters.
\textit{Optimisation} refers to the systematic reduction of cost per unit of
business value through model routing, prompt engineering, infrastructure tier
selection, and caching.

AI FinOps is distinct from, but related to, three adjacent disciplines.
\textit{Cloud FinOps} is the superset from which AI FinOps inherits vocabulary
and organisational structure, but is insufficient for the AI cost structure as
documented in Section~\ref{sec:failure}. \textit{MLOps} addresses the operational
lifecycle of machine learning models --- development, deployment, monitoring --- but
is not primarily concerned with financial governance. \textit{AI Governance} in the
risk and ethics sense addresses issues of model bias, fairness, and accountability,
but does not address cost management. AI FinOps occupies a distinct domain at the
intersection of financial management and AI operations.

\subsection{The AI Cost Taxonomy}
\label{sec:taxonomy}

The AI Cost Taxonomy is the measurement foundation of AI FinOps. It identifies
eight cost categories that collectively account for the total expenditure of a
production AI deployment. For each category, Table~\ref{tab:taxonomy} specifies
the measurable unit, the attribution dimension, and the primary optimisation lever.

\begin{table*}[t]
\centering
\caption{The AI Cost Taxonomy: Eight Dimensions, Measurement Units, and Optimisation Levers}
\label{tab:taxonomy}
\small
\begin{tabularx}{\textwidth}{L{2.0cm} L{1.8cm} L{2.5cm} L{4.8cm}}
\toprule
\textbf{Category} & \textbf{Unit} & \textbf{Attribution} & \textbf{Optimisation Lever} \\
\midrule
Input tokens      & Tokens/request  & Application, team, task type    & Prompt compression; caching; context window management \\
\addlinespace
Output tokens     & Tokens/response & Application, verbosity profile  & Output length constraints; structured response formats \\
\addlinespace
Reasoning tokens  & Tokens/chain    & Model config, task complexity   & Thinking budgets; model routing away from reasoning tiers \\
\addlinespace
Tool call overhead & Calls/session  & Agent type, workflow            & Precision retrieval; tool call batching; result compression \\
\addlinespace
Retry \& failure  & Tokens in failed runs & Infra tier, agent type   & Checkpointing; retry limits; success rate monitoring \\
\addlinespace
GPU compute (training) & GPU-hours/job & Model, training run, team    & Spot instances; job scheduling; checkpoint frequency \\
\addlinespace
GPU compute (inference) & GPU-hours/deployment & Application, utilisation & Right-sizing; batch inference; model quantisation \\
\addlinespace
Storage \& egress & GB stored/transferred & Dataset, model, region   & Data locality; weight caching; egress minimisation \\
\bottomrule
\end{tabularx}
\end{table*}

The first three categories --- input tokens, output tokens, and reasoning tokens ---
constitute 70--90\% of total AI variable cost in most production deployments.
Instrumentation priority should therefore concentrate on these three before extending
to the remaining categories.

Reasoning tokens merit particular emphasis. They are billed at output token rates,
are not surfaced in most AI API response objects by default, and can represent
30--70\% of total token spend in deployments that use extended thinking or
chain-of-thought models indiscriminately. The failure to track reasoning tokens
separately from standard output tokens is the single most common measurement
gap observed in enterprise AI deployments \citep{anthropic2025extended}.

\subsection{The AI FinOps Maturity Model}
\label{sec:maturity}

The AI FinOps Maturity Model defines three levels of organisational capability.
The levels are cumulative: each level incorporates and extends the capabilities
of the levels below it. Figure~\ref{fig:maturity} illustrates the progression;
Table~\ref{tab:maturity} defines the specific capabilities at each level.

\begin{figure}[H]
\centering
\small
\begin{tabular}{ccc}
\toprule
\textbf{Level 1} & \textbf{Level 2} & \textbf{Level 3} \\
\textit{Reactive} & \textit{Managed} & \textit{Optimised} \\
\midrule
Post-hoc         & Proactive        & Automated \\
Bill review      & Alerting         & Optimisation \\
Opaque spend     & Tagged by app    & Cost per outcome \\
\bottomrule
\end{tabular}
\caption{AI FinOps Maturity Model: Three levels of organisational capability.}
\label{fig:maturity}
\end{figure}

\begin{table*}[t]
\centering
\caption{AI FinOps Maturity Model: Capability Definitions by Level}
\label{tab:maturity}
\small
\begin{tabularx}{\textwidth}{L{2.1cm} L{3.8cm} L{3.8cm} L{3.8cm}}
\toprule
\textbf{Capability} & \textbf{Level 1 --- Reactive} & \textbf{Level 2 --- Managed} & \textbf{Level 3 --- Optimised} \\
\midrule
Cost visibility &
  Single line item: ``AI spend'' &
  Cost tagged by application and team &
  Cost attributed by task, model tier, agent, and outcome \\
\addlinespace
Budget governance &
  Post-hoc bill review &
  Spend alerts on monthly thresholds &
  Real-time token budgets with automated circuit breakers \\
\addlinespace
Anomaly detection &
  Finance team notices spike at month end &
  Alert when spend exceeds 120\% of forecast &
  Alert when per-session token rate exceeds 2$\times$ baseline \\
\addlinespace
Model selection &
  Default to one model for all tasks &
  Model selection documented per application &
  Automated routing to optimal tier per query type \\
\addlinespace
Agentic controls &
  No iteration or budget limits &
  Manual review for sessions exceeding cost threshold &
  Automated hard limits; agent surfaces state at budget boundary \\
\addlinespace
Optimisation cadence &
  Ad hoc, when bills surprise &
  Quarterly model and spend review &
  Continuous optimisation with automated recommendations \\
\addlinespace
FinOps ownership &
  Engineering team owns all cost decisions &
  Shared: engineering + finance &
  Dedicated AI FinOps function with tooling and accountability \\
\bottomrule
\end{tabularx}
\end{table*}

Assessment of current enterprise state: the majority of organisations deploying
AI at scale are operating at Level~1. This assessment is consistent with
anonymised patterns observed across enterprise healthcare AI deployments and
with the absence of any widely adopted AI FinOps tooling or standards as of
early 2026. The absence of such standards is itself evidence of the nascent
state of the discipline --- a gap this paper aims to begin closing.

The transition from Level~1 to Level~2 is primarily a technical project:
implementing token logging middleware, deploying a cost attribution dashboard,
and establishing alerting thresholds. For most organisations, this is a
four-to-eight week engineering effort. The transition from Level~2 to Level~3
is primarily an organisational project: establishing the AI FinOps function,
implementing automated routing and budget enforcement, and building the
cost-per-outcome metrics that connect spend to value. This transition requires
six to eighteen months depending on the scale and complexity of the AI portfolio.

\subsection{The Five Governance Questions}
\label{sec:questions}

The Five Governance Questions constitute a minimum viable baseline assessment
instrument for any organisation deploying AI at scale. They are designed to be
answerable by an CAIO or AI platform lead without advanced analytics tooling
--- requiring only that basic token instrumentation and cost attribution are
in place. An organisation that cannot answer all five is, by definition, at
Level~1 maturity.

\begin{enumerate}[leftmargin=*, label=\textbf{Q\arabic*.}, itemsep=4pt]

\item \textbf{What is our total AI spend this month, broken down by application?}\\
\textit{If unanswerable:} AI costs are managed as a black box. Implement application-level
cost tagging on all API calls as the first priority.

\item \textbf{Which application or agent consumed the most tokens last week?}\\
\textit{If unanswerable:} Largest cost drivers cannot be identified. Add per-request
token logging to all AI application middleware.

\item \textbf{What percentage of our token spend is reasoning tokens vs.\ standard output?}\\
\textit{If unanswerable:} The most expensive cost category is invisible. Configure
separate metric tracking for \texttt{thinking\_tokens} vs.\ \texttt{output\_tokens}
in the observability stack.

\item \textbf{Did any agent or application exceed its token budget in the last 30 days?}\\
\textit{If unanswerable:} No budget enforcement exists; discovery is entirely post-hoc.
Implement soft alerts at 80\% of budget and hard stops at 100\% for all production agents.

\item \textbf{What is our cost per successful AI-driven outcome, by use case?}\\
\textit{If unanswerable:} AI spend cannot be connected to business value. Define one
measurable outcome metric per use case and track cost against it monthly.

\end{enumerate}

These questions build in complexity. Questions 1--2 address cost visibility;
question 3 addresses cost-type granularity; question 4 addresses governance
enforcement; question 5 addresses value connection. Organisations should work
through them in sequence, as each question's answerability depends on the
instrumentation required to answer the questions preceding it.

% ─────────────────────────────────────────────────────────────────
\section{Implementation: The 30-Day Visibility Sprint}
\label{sec:implementation}

The 30-day visibility sprint is a structured implementation sequence that delivers
Level~2 maturity baseline capability --- the ability to answer Questions 1 and 2 ---
within a calendar month from a Level~1 starting point. It is designed to be
executable by a small cross-functional team (two to three engineers, one finance
representative) without major tooling purchases or organisational restructuring.

\textbf{Week 1: Instrumentation audit.} Inventory every AI application in production.
For each, document the model tier, monthly token consumption (if tracked), and
whether per-request logging is implemented. Identify the three highest-spend
applications as the priority instrumentation targets.

\textbf{Week 2: Minimum viable logging.} Add the following ten fields to every
production AI application lacking per-request logging: \texttt{timestamp},
\texttt{application\_id}, \texttt{team\_id}, \texttt{model\_id},
\texttt{input\_token\_count}, \texttt{output\_token\_count},
\texttt{reasoning\_token\_count}, \texttt{tool\_calls\_count}, \texttt{session\_id},
\texttt{task\_type}. This is a middleware addition requiring no changes to
application logic or model configuration, typically completable in one to three
engineering days per application.

\textbf{Week 3: Dashboard and alerting.} Aggregate token logs into a basic
reporting view showing total spend by application, model, and team for the
trailing 30 days. Set a spend alert on any application exceeding 150\% of
its rolling 30-day daily average in a single day.

\textbf{Week 4: First governance meeting.} Convene a 60-minute cross-functional
AI cost review. Present 30-day spend data. Identify the top three cost drivers.
Assign budget owners for the three highest-spend applications. Set spend targets
for the following quarter.

\begin{callout}
\small\textbf{Observed outcome:} Across healthcare AI deployments that completed
this sprint, every organisation discovered at least one application consuming
3--5$\times$ its team's believed spend. In the majority of cases, this discovery
resulted in immediate remediation actions --- typically reasoning model opt-out
or agentic loop iteration caps --- that recovered 20--35\% of total AI spend
within the following 30 days.
\end{callout}

% ─────────────────────────────────────────────────────────────────
\section{Empirical Evidence and Validation}
\label{sec:evidence}

\subsection{Market Evidence}

The financial risk of ungoverned AI spend is quantified by several independent
research streams. Gartner projects that by 2030, enterprises failing to optimise
their AI compute environment will pay a 50\% cost premium relative to governed
peers for identical workloads \citep{gartner2025forecast}. This projection is
consistent with the governance dividend observed in practice: organisations that
implement the AI FinOps framework described in this paper typically achieve
40--55\% total cost reduction relative to their pre-governance baseline within
twelve months.

The infrastructure economics literature provides the utilisation-based inflection
points that govern infrastructure tier selection. Uptime Institute (2025) documents
that on-premise GPU deployments become cost-competitive with hyperscaler pricing
at 22\% average utilisation, and with neocloud pricing at 66\% utilisation
\citep{uptime2025}. CoreWeave's 2025 S-1 filing documents revenue of \$1.9 billion
(730\% year-over-year growth) and a \$55.6 billion contracted backlog, providing
market validation of the 40--66\% neocloud cost advantage relative to hyperscaler
GPU pricing \citep{coreweave2025}.

Token price deflation has proceeded at approximately 10$\times$ per year since
2022 across all major model tiers \citep{introl2025}. Midjourney's 2024
migration from hyperscaler GPU infrastructure to Google TPU v6e resulted in a
67\% cost reduction and \$16.8 million in annualised savings \citep{midjourney2024},
providing a validated benchmark for infrastructure tier optimisation at production scale.

\subsection{Observed Enterprise Patterns}

The following patterns are drawn from anonymised assessment data across enterprise
AI deployments in healthcare payer and health IT services environments. No client
names or identifying characteristics are disclosed.

\textbf{Pattern 1: Systematic Level~1 prevalence.} Of organisations assessed
prior to implementing AI FinOps practices, all were operating at Level~1 by the
maturity model defined in Section~\ref{sec:maturity}. None could answer more than
two of the Five Governance Questions (Section~\ref{sec:questions}) without
manual analysis of raw billing records.

\textbf{Pattern 2: Agentic undercounting.} In every deployment where agentic
workloads were present, the engineering team's estimate of agent token consumption
was lower than actual consumption by a factor of 3--10$\times$. The primary
driver in the majority of cases was reasoning token consumption that was neither
tracked nor anticipated in initial cost models.

\textbf{Pattern 3: Governance dividend.} Deployments that implemented token
logging, model routing, and session budget controls --- corresponding to
Level~2 maturity --- achieved cost reductions of 20--55\% relative to
pre-governance baseline within the first 90 days. The lower bound (20\%) was
observed in deployments where agentic workloads were minimal; the upper bound
(55\%) in deployments with significant agentic and reasoning model usage.

\textbf{Pattern 4: Discovery as first-order value.} In every assessed deployment,
the instrumentation sprint (Section~\ref{sec:implementation}) produced at least
one cost anomaly that had been invisible prior to token logging. The act of
measurement itself --- independent of any subsequent optimisation action --- was
consistently cited as the highest-value immediate output of the AI FinOps programme.

% ─────────────────────────────────────────────────────────────────
\section{Discussion}
\label{sec:discussion}

\subsection{The Strategic Dimension of AI FinOps}

The framing of AI FinOps as a cost reduction programme, while accurate, understates
its strategic significance. Organisations that govern AI effectively do not simply
spend less on commodity AI infrastructure --- they redirect the savings toward
the proprietary AI investments that create durable competitive advantage. These
investments take the form of domain-specific knowledge graphs, fine-tuned models
trained on proprietary case data, outcome-linked cost benchmarks, and validated
ontologies built from institutional expertise. These assets cannot be replicated
without years of the same case history and governance infrastructure.

This creates a compounding dynamic. Better governance reduces spend on commodity
AI. Reduced spend on commodity AI frees capital for proprietary AI investment.
Proprietary AI investment improves outcomes, which generates more case data,
which matures the knowledge assets further. Organisations that fail to initiate
this cycle remain in a pattern of increasing spend on ungoverned commodity AI
with no proprietary differentiation to show for it.

\subsection{Implications for the FinOps Foundation}

The FinOps Foundation's existing frameworks represent a strong foundation for
extending into AI. Three specific extensions are recommended. First, the FinOps
persona framework should be extended to include the AI FinOps Analyst, AI Cost
Engineer, and AI Budget Owner roles defined in this paper, with capability
profiles analogous to those established for cloud FinOps personas. Second, the
FinOps capability taxonomy should be extended with AI-specific capability
domains: token cost attribution, agentic budget governance, model portfolio
management, and AI ROI measurement. Third, the Five Governance Questions
proposed in Section~\ref{sec:questions} should be incorporated into the
FinOps Foundation's maturity assessment tooling as a validated AI FinOps
baseline instrument.

\subsection{Limitations}

This paper has three primary limitations. First, the empirical patterns in
Section~\ref{sec:evidence} are drawn from a single industry sector (healthcare
payer and health IT services) and may not generalise uniformly to other sectors
with different AI workload profiles. Second, the specific pricing benchmarks
referenced throughout reflect market conditions as of Q1 2026; given the 10$\times$
annual token price deflation documented in Section~\ref{sec:background}, these
figures will become dated within 12--18 months, though the governance framework
itself is designed to be robust to price changes. Third, the maturity model
and governance questions are practitioner-validated rather than validated
through formal survey methods; future work should develop and administer a
psychometrically validated instrument.

\subsection{Future Research Directions}

Four directions for future research are identified. First, a formal survey
study of AI FinOps maturity across enterprise sectors would validate the
prevalence of Level~1 and enable regression analysis of governance
capability drivers. Second, longitudinal study of the governance dividend ---
tracking cost reduction as organisations progress through maturity levels
over 12--24 months --- would provide causal evidence for the framework's
economic claims. Third, the Five Governance Questions should be developed
into a psychometrically validated assessment instrument with established
reliability and construct validity. Fourth, the extension of AI FinOps
frameworks to multi-agent and cross-organisational AI workflows --- where
cost attribution must span organisational boundaries --- represents an
important frontier as agent-to-agent collaboration matures.

% ─────────────────────────────────────────────────────────────────
\section{Conclusion}
\label{sec:conclusion}

Enterprise AI is accumulating financial risk faster than enterprise governance
is accumulating the capability to manage it. This paper has argued that the
primary cause of this gap is not organisational negligence but architectural
mismatch: the Cloud FinOps frameworks that organisations rely on were built
for a cost structure --- compute hours, resource-level attribution, linear
scaling --- that AI systems do not have. AI systems are billed in tokens,
attributed at the request level, and capable of nonlinear cost dynamics
through reasoning chains and agentic loops that Cloud FinOps cannot detect,
attribute, or bound.

We have introduced four constructs to address this gap. The \textit{definition
of AI FinOps} establishes the discipline as a formal organisational capability
distinct from Cloud FinOps, MLOps, and AI risk governance. The \textit{AI Cost
Taxonomy} provides the eight-dimension measurement foundation without which
no governance action is possible. The \textit{AI FinOps Maturity Model} gives
organisations a progression path from reactive bill review to automated
optimisation, with specific, assessable capabilities at each level. The
\textit{Five Governance Questions} constitute a minimum viable baseline
instrument that any CAIO can apply today, without advanced tooling, to
determine whether their organisation's AI spend is under governance.

The Jevons Paradox ensures that these governance capabilities will not become
less important as token prices fall. Every order-of-magnitude decline in token
cost enables new classes of AI application that expand total consumption, creating
new governance challenges at larger scale. The organisations that build AI FinOps
capability now --- before their AI spend reaches the scale at which ungoverned
consumption becomes a crisis --- will outperform those that build it reactively.
The discipline will be defined. The question is whether practitioners define it
proactively, or wait to be defined by the consequences of not doing so.

% ─────────────────────────────────────────────────────────────────
\balance
\bibliographystyle{plainnat}
\bibliography{aifinops}

% ─────────────────────────────────────────────────────────────────
\onecolumn
\appendix

\section*{Appendix A: AI FinOps Capability Assessment Checklist}
\label{app:checklist}
\addcontentsline{toc}{section}{Appendix A: Capability Assessment Checklist}

The following 14-item checklist operationalises the AI FinOps Maturity Model
for self-assessment by CAIOs and AI Platform Leads. Mark each capability as
\textbf{In Place} ($\checkmark$) or \textbf{Not in Place} ($\square$). Count
of $\checkmark$ marks: 0--4 = Level~1; 5--10 = Level~2 transition; 11--14 = Level~2/3.

\vspace{6pt}
\begin{tabular}{lp{9cm}cc}
\toprule
\textbf{\#} & \textbf{Capability} & \textbf{Not in Place} & \textbf{In Place} \\
\midrule
1  & Token logging on all production AI applications              & $\square$ & $\square$ \\
2  & Cost attribution by application and team                     & $\square$ & $\square$ \\
3  & Model-task fit documented per application                    & $\square$ & $\square$ \\
4  & Reasoning token tracking separate from output tokens         & $\square$ & $\square$ \\
5  & Hard session budgets on all production agents                & $\square$ & $\square$ \\
6  & Iteration limits on all agentic loops                        & $\square$ & $\square$ \\
7  & Prompt caching on high-volume system prompts                 & $\square$ & $\square$ \\
8  & GPU utilisation measured by workload                         & $\square$ & $\square$ \\
9  & Infrastructure tier matched to utilisation profile           & $\square$ & $\square$ \\
10 & Quarterly model evaluation cadence established               & $\square$ & $\square$ \\
11 & Open-source model evaluation for eligible workloads          & $\square$ & $\square$ \\
12 & AI budget structured with scenario range (not single figure) & $\square$ & $\square$ \\
13 & ROI bridge completed for every active AI use case            & $\square$ & $\square$ \\
14 & Quarterly AI cost review cadence established                 & $\square$ & $\square$ \\
\bottomrule
\end{tabular}

\vspace{12pt}
\section*{Appendix B: Minimum Viable Token Logging Specification}
\label{app:logging}
\addcontentsline{toc}{section}{Appendix B: Minimum Viable Token Logging Specification}

The following ten-field schema constitutes the minimum viable log record for AI
FinOps instrumentation. It should be emitted by application middleware on every
AI API call and stored in a queryable log store (e.g., BigQuery, Snowflake,
Elasticsearch). All fields are required; missing fields should be logged as
\texttt{null} rather than omitted to preserve schema consistency.

\vspace{6pt}
\begin{tabular}{llp{8cm}}
\toprule
\textbf{Field} & \textbf{Type} & \textbf{Description} \\
\midrule
\texttt{timestamp}              & ISO 8601 datetime & UTC time of API call completion \\
\texttt{application\_id}        & string            & Unique identifier of the calling application \\
\texttt{team\_id}               & string            & Team or cost centre owning the application \\
\texttt{model\_id}              & string            & Model identifier as returned by the API \\
\texttt{input\_token\_count}    & integer           & Tokens in the full request context \\
\texttt{output\_token\_count}   & integer           & Tokens in the model's visible response \\
\texttt{reasoning\_token\_count}& integer           & Thinking/reasoning tokens if applicable; \texttt{null} otherwise \\
\texttt{tool\_calls\_count}     & integer           & Number of tool invocations in session; \texttt{null} for non-agentic \\
\texttt{session\_id}            & string            & Shared identifier across all calls in a multi-turn session \\
\texttt{task\_type}             & enum              & Controlled vocabulary: \texttt{classification}, \texttt{summarisation}, \texttt{generation}, \texttt{retrieval}, \texttt{reasoning}, \texttt{agentic}, \texttt{other} \\
\bottomrule
\end{tabular}

\end{document}
